% Это основная команда, с которой начинается любой \LaTeX-файл. Она отвечает за тип документа, с которым связаны основные правил оформления текста.
\documentclass[14pt]{extarticle}

% Здесь идет преамбула документа, тут пишутся команды, которые настраивают LaTeX окружение, подключаете внешние пакеты, определяете свои команды и окружения. В данном случае я это делаю в отдельных файлах, а тут подключаю эти файлы.

% Здесь я подключаю разные стилевые пакеты. Например возможности набирать особые символы или возможность компилировать русский текст. Подробное описание внутри.
\usepackage{packages}

% Здесь я определяю разные окружения, например, теоремы, определения, замечания и так далее. У этих окружений разные стили оформления, кроме того, эти окружения могут быть нумерованными или нет. Все подробно объяснено внутри.
\usepackage{environments}

% Здесь я определяю разные команды, которых нет в LaTeX, но мне нужны, например, команда \tr для обозначения следа матрицы. Или я переопределяю LaTeX команды, которые работают не так, как мне хотелось бы. Типичный пример мнимая и вещественная часть комплексного числа \Im, \Re. В оригинале они выглядят не так, как мы привыкли. Кроме того, \Im еще используется и для обозначения образа линейного отображения. Подробнее описано внутри.
\usepackage{commands}

% fix table of contents
\usepackage[numbib,notlof,notlot,nottoc]{tocbibind}

\counterwithin{figure}{section}

% \renewcommand{\thefigure}{\thesection.\arabic{figure}}

% С этого момента начинается текст документа
\begin{document}

\counterwithin{lstlisting}{section}

\thispagestyle{empty}
титульный лист
\pagebreak

\thispagestyle{empty}

\selectlanguage{russian}
\begin{abstract}

Аннотация на русском языке

\end{abstract}
\pagebreak

\thispagestyle{empty}

\selectlanguage{english}
\begin{abstract}

Abstract in english.

\end{abstract}
\pagebreak

\selectlanguage{russian}

% Здесь будет автоматически генерироваться содержание документа
{
\hypersetup{linkcolor=black}
\tableofcontents
}

% Каждая секция должна начинаться с новой страницы
\pagebreak
\section{Введение}

Пример цитирования ~\cite{GabRam}.

Пример ссылки на секцию.

В сеции \ref{part2} будут описаны существующие работы и решения.

\pagebreak
\section{Существующие работы и решения}\label{part2}

\pagebreak
\section{Описание реализованных алгоритмов}\label{part3}

\pagebreak
\section{Детали реализации}\label{part4}

\pagebreak
\section{Тестирование корректности и производительности}\label{part5}

Пример кода
\begin{lstlisting}[language=bash, label=delay,caption={Добавление задержки}]
  qdisc add dev eth0 root netem $DELAY
\end{lstlisting}

Пример добавления svg картинок

\begin{figure}[H]
  \centering
    \subfloat[\centering Задержка записей]{{\includesvg[width=8cm]{write_3} }}%
    \qquad
    \subfloat[\centering Задержка чтений]{{\includesvg[width=8cm]{read_3} }}%
    \caption{Зависимость задержки записей и чтений от задержки между контейнерами}
\end{figure}

\pagebreak
\section{Заключение}\label{part6}

% Здесь автоматически генерируется библиография. Первая команда задает стиль оформления библиографии, а вторая указывает на имя файла с расширением bib, в котором находится информация об источниках. Чтобы источник появился, надо на него сослаться.
\pagebreak
\bibliographystyle{ugost2008}
\renewcommand{\refname}{Библиографический список}
\bibliography{bibl}


% С этого момента глобальная нумерация идет буквами. Этот раздел я добавил лишь для демонстрации возможностей LaTeX, его можно и нужно удалить и он не нужен для курсового проекта непосредственно.
% \appendix
%
% \pagebreak
% Проведем небольшой обзор возможностей \LaTeX. Далее идет % обзорный кусок, который надо будет вырезать. Он приведен % лишь для демонстрации возможностей \LaTeX.
%
% \section{Нумеруемый заголовок}
% Текст раздела
% \subsection{Нумеруемый подзаголовок}
% Текст подраздела
% \subsubsection{Нумеруемый подподзаголовок}
% Текст подподраздела
%
% \section*{Не нумеруемый заголовок}
% Текст раздела
% \subsection*{Не нумеруемый подзаголовок}
% Текст подраздела
% \subsubsection*{Не нумеруемый подподзаголовок}
% Текст подподраздела
%
%
% \paragraph{Заголовок абзаца} Текст абзаца
%
% Формулы в тексте набирают так $x = e^{\pi i}\sqrt{\text% {формула}}$. Выключенные не нумерованные формулы набираются % либо так:
% \[
% x = e^{\pi i}\sqrt{\text{формула}}
% \]
% Либо так
% $$
% x = e^{\pi i}\sqrt{\text{формула}}
% $$
% Первый способ предпочтительнее при подаче статей в журналы % AMS, потому рекомендую привыкать к нему.
%
% Выключенные нумерованные формулы:
% \begin{equation}\label{Equation1}
% % \label{имя-метки} эта команда ставит метку, на которую % потом можно сослаться с помощью \ref{имя-метки}. Метки можно % ставить на все объекты, у которых есть автоматические % счетчики (номера разделов, подразделов, теорем, лемм, формул % и т.д.
% x = e^{\pi i}\sqrt{\text{формула}}
% \end{equation}
% Или не нумерованная версия
% \begin{equation*}
% x = e^{\pi i}\sqrt{\text{формула}}
% \end{equation*}
%
% Уравнение~\ref{Equation1} радостно занумеровано.
%
% Лесенка для длинных формул
% \begin{multline}
% x = e^{\pi i}\sqrt{\text{очень очень очень длинная формула}}% =\\
% \tr A - \sin(\text{еще одна очень очень длинная формула})=\\
% \cos z \Im \varphi(\text{и последняя длинная при длинная % формула})
% \end{multline}
%
% Многострочная формула с центровкой
% \begin{gather}
% x = e^{\pi i}\sqrt{\text{очень очень очень длинная формула}}% =\\
% \tr A - \sin(\text{еще одна очень очень длинная формула})=\\
% \cos z \Im \varphi(\text{и последняя длинная при длинная % формула})
% \end{gather}
%
% Многострочная формула с ручным выравниванием. Выравнивание % идет по знаку $\&$, который на печать не выводится.
% \begin{align}
% x = &e^{\pi i}\sqrt{\text{очень очень очень длинная формула}}% =\\
% &\tr A - \sin(\text{еще одна очень очень длинная формула})=\\
% &\cos z \Im \varphi(\text{и последняя длинная при длинная % формула})
% \end{align}
%
% \begin{theorem}
% Текст теоремы
% \end{theorem}
% \begin{proof}
% В специальном окружении оформляется доказательство.
% \end{proof}
%
% \begin{theorem}[Имя теоремы]
% Текст теоремы
% \end{theorem}
% \begin{proof}[Доказательство нашей теоремы]
% В специальном окружении оформляется доказательство.
% \end{proof}
%
% \begin{definition}
% Текст определения
% \end{definition}
%
% \begin{remark}
% Текст замечания
% \end{remark}
%
% \paragraph{Перечни:} Нумерованные
% \begin{enumerate}
% \item Первый
% \item Второй
% \begin{enumerate}
% \item Вложенный первый
% \item Вложенный второй
% \end{enumerate}
% \end{enumerate}
%
% Не нумерованные
%
% \begin{itemize}
% \item Первый
% \item Второй
% \begin{itemize}
% \item Вложенный первый
% \item Вложенный второй
% \end{itemize}
% \end{itemize}


% Здесь текст документа заканчивается
\end{document}
% Начиная с этого момента весь текст LaTeX игнорирует, можете вставлять любую абракадабру.
